\documentclass[11pt,a4paper]{article}
\usepackage[T1]{fontenc}
\usepackage[utf8]{inputenc}
\usepackage{lmodern,microtype}
\usepackage[a4paper,margin=27mm,headheight=14pt]{geometry}
\usepackage{amsmath,amssymb,amsthm,mathtools,bm}
\usepackage{booktabs,array,longtable,enumitem}
\usepackage[dvipsnames]{xcolor}
\usepackage{hyperref}
\hypersetup{colorlinks=true,linkcolor=MidnightBlue,citecolor=MidnightBlue,urlcolor=MidnightBlue,
 pdftitle={Signed Condensation in Quadratic-Feedback Transseries},
 pdfauthor={Research report prepared for Vladimir Reshetnikov}}
\usepackage{fancyhdr}
\pagestyle{fancy}\fancyhf{}
\fancyhead[L]{\small Signed condensation and sharp reversion}
\fancyhead[R]{\small September 2026}\fancyfoot[C]{\thepage}
\setlength{\emergencystretch}{2em}
\numberwithin{equation}{section}
\newtheorem{theorem}{Theorem}[section]
\newtheorem{lemma}[theorem]{Lemma}
\newtheorem{proposition}[theorem]{Proposition}
\newtheorem{corollary}[theorem]{Corollary}
\theoremstyle{definition}\newtheorem{definition}[theorem]{Definition}
\newtheorem{question}{Research question}
\theoremstyle{remark}\newtheorem{remark}[theorem]{Remark}
% ed. (2026-09-29): unnumbered environment for ProveIt editorial notes; it does not shift any numbering.
\theoremstyle{definition}
\newtheorem*{ednoteinner}{Editorial note (ProveIt, 2026-09-29)}
\newenvironment{ednote}{\begin{ednoteinner}\sloppy\Urlmuskip=0mu plus 1mu\relax}{\end{ednoteinner}}
\newcommand{\R}{\mathbb R}
\newcommand{\N}{\mathbb N}
\newcommand{\U}{\widehat U}
\newcommand{\Q}{\widehat Q}
\newcommand{\Bcal}{\mathcal B}
\newcommand{\dd}{\,\mathrm d}
\newcommand{\ee}{\mathrm e}
\newcommand{\coef}[2]{[#1^{#2}]}
\newcommand{\abs}[1]{\left|#1\right|}
\newcommand{\RepoCommit}{aab173a7ccdd48add72191d68ce6b0c12ffd72fe}
\begin{document}
% ed. (2026-09-29): no page anchors on the title page, which the titlepage environment numbers 1 like the following page (removes a duplicate-destination warning).
\hypersetup{pageanchor=false}
\begin{titlepage}
\centering
\vspace*{15mm}
{\large PROVEIT RESEARCH CONTINUATION\par}
\vspace{15mm}
{\LARGE\bfseries Signed Condensation in\\[3mm]
Quadratic-Feedback Transseries\par}
\vspace{7mm}
{\Large Sharp reversion, universal cancellation,\\
and inverse Borel growth\par}
\vspace{14mm}
{\large Research report prepared for Vladimir Reshetnikov\par}
\vspace{4mm}
{September 29, 2026\par}
\vspace{15mm}
\begin{minipage}{0.91\textwidth}
\begin{abstract}
For the formal solution of
\(\U(q)=\sum_{j\ge1}w_jq^j\exp(\lambda_j\U(q))\), suppose that
\(w_1=1\), all primitive weights and slopes are nonnegative, and, outside a
finite prefix, \(w_j=1\) and \(\lambda_j=aj^2+bj+d\), with \(a>0\).
We determine a full multiplicative equivalent for the coefficients of its
compositional inverse. With \(\beta=\lambda_1+w_2\),
\(r_n(1+r_n)e^{2r_n}=an\), and
\(S_n=\exp(nr_n(1+2r_n)/(1+r_n))/\sqrt{1+4r_n+2r_n^2}\), the result is
\[
 [u^n]\Q(u)\sim-S_n\exp\!\left(\frac{b-\beta}{a}r_n\right).
\]
For \(w_j=1\) and \(\lambda_j=aj^2\), this proves the precise quadratic
inverse conjecture in the repository's finite-core article. The essential
step is a signed-tail localization theorem: after a sufficiently large but
fixed analytic core is removed, configurations containing at least two
remaining actions contribute less than \(S_n n^{-A}\), for any prescribed
\(A>0\). It is proved at the level of absolute configuration weights,
before cancellations are used. We then derive eventual negativity, a
finite-prefix sensitivity law, the exact positive-axis growth equivalent of
the inverse Borel transform, and the least inverse-term equivalent. The
proofs are self-contained apart from classical inversion and elementary
asymptotic tools.
An exact verification program, recorded integer coefficients, a source
audit, and a research agenda accompany the manuscript.
\end{abstract}
\end{minipage}
\vfill
\begin{minipage}{0.91\textwidth}\small
\textbf{Status and scope.} This is a research manuscript with proofs for
independent review, not a Lean-verified development. The novelty claim is
scoped to the explicitly audited repository conjecture and the extensions
proved here; no claim of resolving general transseries summability or a
classical field-wide conjecture is made.
\end{minipage}
\end{titlepage}
\hypersetup{pageanchor=true}
\tableofcontents
\newpage

\section{The specific unresolved problem}
\label{sec:context}

\subsection{Repository provenance}
The repository \emph{ProveIt} separates its formal transseries library from
its mathematical research reports. At the inspected snapshot, the project
README expressly records that the recent research arrivals have not been
reviewed claim by claim and do not become Lean theorems merely by being
filed alongside the library \cite{repo,inventory}. We retain that distinction.
The fixed snapshot used throughout this report is
\begin{center}\small\texttt{\RepoCommit}.\end{center}
The principal source is
\begin{quote}\small\raggedright
\path{Analysis/Transseries/docs/series-and-transseries/}\allowbreak
\path{Finite_Core_Universality_Exponential_Feedback/article.tex}.
\end{quote}
Its blob identifier is
\begin{center}\small\texttt{c2b6faf623003111e9d9e6bbb162a67429f3ce42}.\end{center}
The source's Section ``Quadratic and subquadratic inverse asymptotics,''
particularly labels \texttt{prop:inverse-response},
\texttt{conj:quadratic-inverse}, and \texttt{eq:inverse-conjecture}, states
an exact first-order inverse response but leaves the cancellation estimate
needed to pass from that response to the full inverse unproved \cite{finitecore}.
The relevant inspected source range is lines 1623--1678; the precise
conjecture is at lines 1661--1668, with its explanation immediately after it.
% ed. (2026-09-29): the cited line numbers predate the editorial pass of 2026-09-29.
\begin{ednote}
The cited lines are those of the blob and snapshot identified above. After the editorial pass of 2026-09-29 the range
1623--1678 is 1666--1721, and the conjecture is at lines 1704--1711 of
\path{Finite_Core_Universality_Exponential_Feedback/article.tex}, followed
by an editorial note recording the claimed proofs, now including this one.
\end{ednote}

For the pure quadratic model, that conjecture is
\begin{equation}
 v_n\sim-S_n e^{-(1+1/a)r_n},
 \qquad \Q(u)=\sum_{n\ge1}v_nu^n.
 \label{eq:oldconjecture}
\end{equation}
The same article proves a forward equivalent
\(u_n\sim S_ne^{(1+1/a)r_n^2}\), but proves
\(v_n\sim-u_n\) only in its superquadratic regime. Consequently,
Gevrey-class preservation, preservation of weighted exponential type, or
an estimate of \(\log|v_n|\) to accuracy \(o(n)\) does not settle
\eqref{eq:oldconjecture}. In the quadratic case the missing factor in the
forward-to-inverse comparison is already exponential in \((\log n)^2\).

\subsection{What is added here}
The principal contribution is a proof of \eqref{eq:oldconjecture}, including
its sign and multiplicative constant. We prove it in the larger class with
an arbitrary nonnegative finite prefix and an eventually quadratic polynomial
tail. The argument has two distinct components. A positive majorant proves
that, after a sufficiently large fixed cutoff, exactly one remaining action
accounts for the inverse coefficient up to an arbitrarily small polynomial
error relative to \(S_n\). A separate analytic-core calculation then sums
the cancellations within that one-action contribution.

This separation matters. Bounding all inverse-composition terms absolutely
at the start and comparing them directly with the desired answer loses the
factor we are trying to determine. Conversely, computing a compelling
one-action saddle does not exclude a larger contribution from the rest of
the signed expansion. Both issues must be resolved.

Two further consequences contain more information than coefficient type:
a multiplicative equivalent for the inverse Borel transform on the positive
real axis, and a multiplicative equivalent for its smallest formal term.
The latter is \emph{not} asserted to be an analytic remainder estimate.

\subsection{Relation to the literature}
Lagrange inversion and its multivariate forms are classical; Gessel's survey
is a convenient reference \cite{gessel}. General inversion with infinitely
many colours is developed by Jansen, Kuna, and Tsagkarogiannis \cite{jkt}.
Our finite coefficient formula is a specialization of that classical
machinery, not a new inversion formula.

Borinsky develops an algebraic calculus for coefficients possessing
fixed factorial asymptotic expansions and proves composition and inversion
rules in that setting \cite{borinsky}. The present sequences have zero type
relative to every fixed exponential multiple of \(n!\), while still having
zero ordinary radius of convergence. Thus a fixed-factorial expansion with
all asymptotic coefficients zero cannot distinguish the cancellation laws
proved below. This is a reason to keep the slowly varying saddle coordinate
as part of the asymptotic data. The general formal transseries viewpoint is
explained by Edgar \cite{edgar}; our support is only \(\{e^{-nx}:n\ge1\}\).
The difficulty lies in large-order coefficients, not in a new support order.

\section{Model and principal theorems}
\label{sec:main}

Fix real numbers \(a>0\), \(b,d\in\R\), and an integer \(J_0\ge2\).
Assume
\begin{equation}
 w_1=1,\quad w_j\ge0,\quad\lambda_j\ge0,\qquad
 w_j=1,\quad\lambda_j=aj^2+bj+d\quad(j>J_0).
 \label{eq:assumptions}
\end{equation}
All exceptional weights and slopes are finite. In particular, a weight in
the finite prefix may vanish. Define
\begin{equation}
 \U(q)=\sum_{j\ge1}w_jq^j e^{\lambda_j\U(q)},\qquad
 \Q=\U^{\langle-1\rangle},\qquad
 \beta=\lambda_1+w_2.
 \label{eq:model}
\end{equation}
These equations are formal in \(\R[[q]]\). The coefficient of degree \(n\)
in the defining equation only uses \(j\le n\) and earlier coefficients of
\(\U\). Hence \(\U\) exists uniquely, \(u_1=1\), and its formal inverse
exists uniquely. Direct expansion gives
\begin{equation}
 u_2=\beta,\qquad v_2=-\beta.
 \label{eq:beta}
\end{equation}

For real \(n>0\), let \(r=r(n)>0\) solve
\begin{equation}
 r(1+r)e^{2r}=an.
 \label{eq:r}
\end{equation}
The left side is strictly increasing from zero to infinity. Put
\begin{equation}
 j_n=\frac{n}{1+r_n},\qquad k_n=\frac{nr_n}{1+r_n},\qquad
 D_n=1+4r_n+2r_n^2,
 \qquad S_n=\frac{e^{k_n(1+2r_n)}}{\sqrt{D_n}}.
 \label{eq:S}
\end{equation}
Here \(j_n\) is a real saddle coordinate, not an integer index. We have
\(r_n\sim\tfrac12\log n\), \(j_n\sim2n/\log n\), and \(k_n\sim n\).
The exact coordinate \(r_n\), rather than a truncated logarithmic expansion,
is essential in a multiplicative equivalent.

\begin{theorem}[Sharp quadratic reversion]
\label{thm:main}
Under \eqref{eq:assumptions},
\begin{equation}
 \boxed{\displaystyle
 v_n\sim-S_n\exp\!\left(\frac{b-\beta}{a}r_n\right).}
 \label{eq:main}
\end{equation}
In particular, \(v_n<0\) for every sufficiently large integer \(n\).
The constant term \(d\) of the tail affects neither the displayed exponential
factor nor its leading multiplicative constant.
\end{theorem}

\begin{corollary}[The repository conjecture]
\label{cor:pure}
For \(w_j=1\), \(\lambda_j=aj^2\) at every \(j\ge1\),
\[
 v_n\sim-\frac{\exp\!\left(
 \frac{nr_n(1+2r_n)}{1+r_n}-(1+1/a)r_n\right)}
 {\sqrt{1+4r_n+2r_n^2}}.
\]
Thus \eqref{eq:oldconjecture} holds for every fixed \(a>0\).
\end{corollary}

\begin{corollary}[Finite-prefix sensitivity]
\label{cor:prefix}
Suppose two models have the same eventual \(a,b,d\) and respective values
\(\beta,\widetilde\beta\). Then
\begin{equation}
 \frac{v_n}{\widetilde v_n}
 \sim\exp\!\left(-\frac{\beta-\widetilde\beta}{a}r_n\right).
 \label{eq:prefix}
\end{equation}
In particular, preserving \(u_2\) makes every other finite-prefix alteration
invisible in the leading inverse equivalent.
\end{corollary}

% ed. (2026-09-29): two earlier-filed independent claimed proofs of the same theorem (batch 49), not visible from this article's snapshot.
\begin{ednote}
Two earlier-filed packages, which this article could not see, claim the
same inverse law, so this article is a third independent claimed proof of
the conjecture \texttt{conj:quadratic-inverse}.
\path{Quadratic_Exponential_Feedback_After_Reversion/} \cite{ed:qfr}
(Theorem~\texttt{thm:main}) has the same hypotheses, with \(b,d,\beta\)
written \(d,e,\mu\); its Corollary~\texttt{cor:universality} is
Corollary~\ref{cor:prefix} here and its Theorem~\texttt{thm:borel} is
Theorem~\ref{thm:borel} here.
\path{Signed_Quadratic_Feedback_Inversion/} \cite{ed:sqi}
(Theorem~\texttt{thm:main}) has an exact tail \(aj^2\), finitely many
exceptional weights and slopes that may be arbitrary real numbers, and the
factor
\(\exp(-(\lambda_1+w_2)r_n/a)\). The three sets of exact coefficients agree
at every common degree (through \(240\) for \(a=1\) and \(160\) for
\(a=2\)). Theorem~\ref{thm:least} sharpens the least-term index corollary of
the second package and adds the value; the localization proof here is a
third route, closest to the first package's positive majorant. None of the
three proofs has been independently reviewed.
\end{ednote}

For the remaining statements, write \(\gamma=(b-\beta)/a\) and use the
ordinary factorial-normalized Borel transform
\begin{equation}
 \Bcal\Q(z)=\sum_{n\ge1}\frac{v_n}{n!}z^n.
 \label{eq:borel-def}
\end{equation}

\begin{theorem}[Inverse Borel growth]
\label{thm:borel}
The series \eqref{eq:borel-def} is entire. For positive real \(z\to\infty\),
let \(R=R(z)>0\) be defined by \(R(R+1)=az\). Then
\begin{equation}
 \boxed{\displaystyle
 \Bcal\Q(z)\sim-
 \frac{e^{\gamma R}}{\sqrt{2R+1}}
 \exp\!\left(\frac{R}{a}e^{2R}\right).}
 \label{eq:borel}
\end{equation}
Consequently its ordinary Borel--Laplace integral in the positive direction
does not exist, for any positive Laplace parameter.
\end{theorem}

\begin{theorem}[Least inverse term]
\label{thm:least}
Choose \(N\) so large that \(v_n<0\) for \(n\ge N\), and set
\(R=\tfrac12\log(1/t)\) for \(t\downarrow0\). Then
\begin{equation}
 \boxed{\displaystyle
 \min_{n\ge N}|v_n|t^n
 \sim\frac{\exp\!\left(-\frac{R^2}{at}+\gamma R\right)}
 {\sqrt{1+4R+2R^2}}.}
 \label{eq:least}
\end{equation}
Every minimizing index satisfies
\begin{equation}
 n_{\min}=\frac{R(R+1)}{at}
       +o\!\left(\sqrt{\frac{R(R+1)}{at}}\right).
 \label{eq:minindex}
\end{equation}
No claim about a uniquely minimizing integer, an \(O(1)\) index error, or an
analytic truncation remainder is included.
\end{theorem}

\section{An exact signed configuration formula}
\label{sec:exact}

The inverse can be defined directly, without first computing \(\U\):
\begin{equation}
 u=\sum_{j\ge1}w_j\Q(u)^j e^{\lambda_j u}.
 \label{eq:inverse-kernel}
\end{equation}
Put \(\Q(u)=u e^{-\lambda_1u}Z(u)\). Then
\begin{equation}
 1=Z+\sum_{j\ge2}x_jZ^j,
 \qquad x_j=w_ju^{j-1}e^{(\lambda_j-j\lambda_1)u}.
 \label{eq:Z}
\end{equation}

\begin{lemma}[Finite coefficient formula]
\label{lem:exact}
For a finitely supported multiplicity vector \(\bm m=(m_j)_{j\ge2}\), write
\begin{equation}
 K=\sum_{j\ge2}m_j,\quad L=\sum_{j\ge2}(j-1)m_j,\quad
 h=n-1-L,\quad E(\bm m)=\sum_{j\ge2}(\lambda_j-j\lambda_1)m_j-\lambda_1.
 \label{eq:config}
\end{equation}
Then
\begin{equation}
 v_n=\frac{(-\lambda_1)^{n-1}}{(n-1)!}
 +\sum_{\substack{\bm m\ne0\\L\le n-1}}
 (-1)^K\frac{(L+K)!}{(L+1)!\prod_jm_j!}
 \left(\prod_jw_j^{m_j}\right)\frac{E(\bm m)^h}{h!}.
 \label{eq:exact}
\end{equation}
The sum is finite. We use \(0^0=1\), as required by the exponential series.
\end{lemma}
\begin{proof}
Set \(Y=1-Z\) in \eqref{eq:Z}, and temporarily replace every \(x_j\) by
\(sx_j\). The equation is \(Y=s\sum_jx_j(1-Y)^j\). Ordinary Lagrange
inversion in \(s\) gives
\[
 [s^K]Y=\frac1K[t^{K-1}]\left(\sum_jx_j(1-t)^j\right)^K.
\]
For a specified multiplicity vector with \(K\ge1\), the multinomial factor
is \(K!/\prod m_j!\), and the exponent of \(1-t\) is
\(\sum jm_j=L+K\). Thus its coefficient in \(Z=1-Y\) is
\[
 -\frac{(K-1)!}{\prod m_j!}(-1)^{K-1}
   \binom{L+K}{K-1}
 =(-1)^K\frac{(L+K)!}{(L+1)!\prod m_j!}.
\]
Substitute the \(x_j\), multiply by \(ue^{-\lambda_1u}\), and take the
coefficient of \(u^n\). This yields \eqref{eq:exact}. The cost \(j-1\ge1\)
forces \(j\le n\), \(K\le L\le n-1\), proving finiteness.
\end{proof}

For the pure model, the exponent simplifies to
\begin{equation}
 E(\bm m)=a\left(\sum_{j\ge2}j(j-1)m_j-1\right)>0
 \qquad(\bm m\ne0).
 \label{eq:pureE}
\end{equation}
Therefore the signs in its nonempty configuration sum are exactly
\((-1)^K\). The initial term in \eqref{eq:exact} is kept separate.
For a general finite prefix, \(E\) can have either sign at small degree;
the localization proof below uses its absolute value until a large action
makes it positive.

\begin{definition}[Absolute configuration mass]
Let \(A_{n,L}\) be the sum of the absolute values of the nonempty terms
in \eqref{eq:exact} with the specified excess \(L\). Let \(A_n\) be their
sum over \(L\). For an integer \(M\ge J_0\), let \(A_{n,M}^{(2)}\) be the
same absolute mass restricted to configurations having at least two
primitive actions \(j>M\), counted with multiplicity.
\end{definition}

These masses are not the absolute value of an already cancelled coefficient.
They are larger, positive quantities. The next two sections prove estimates
for them strong enough to justify later signed manipulations.

\section{The bare saddle and its weighted sum}
\label{sec:bare}

Define, for integers \(1\le j\le n\),
\begin{equation}
 T_{n,j}=\frac{(aj^2)^{n-j}}{(n-j)!},\qquad
 F_n(x)=(n-x)\left(1+\log\frac{ax^2}{n-x}\right)
 \quad(0<x<n).
 \label{eq:T}
\end{equation}
Differentiation gives
\begin{align}
 F_n'(x)&=-\log\frac{ax^2}{n-x}+\frac{2(n-x)}x,\label{eq:Fprime}\\
 F_n''(x)&=-\frac4x-\frac1{n-x}-\frac{2(n-x)}{x^2}<0.
 \label{eq:Fsecond}
\end{align}
Thus the real saddle is unique. Equations \eqref{eq:r}--\eqref{eq:S} give
\begin{equation}
 F_n'(j_n)=0,\quad F_n(j_n)=k_n(1+2r_n),\quad
 -F_n''(j_n)=\frac{D_n}{k_n},\quad
 \sigma_n^2=\frac{k_n}{D_n}\asymp\frac{n}{(\log n)^2}.
 \label{eq:baresaddle}
\end{equation}

\begin{lemma}[Bare and slowly weighted saddle sums]
\label{lem:bare}
For every fixed real \(c\),
\begin{equation}
 \sum_{j=1}^nT_{n,j}\exp\!\left(c\frac{n-j}{j}\right)
 \sim S_ne^{cr_n}.
 \label{eq:weightedbare}
\end{equation}
The sum may be restricted to \(n/\log n\le j\le3n/\log n\).
The discarded part is exponentially small relative to the full sum.
Also, in a fixed sufficiently small relative neighborhood of \(j_n\),
\begin{equation}
 F_n(j_n)-F_n(j)\ge c_0\frac{(\log n)^2}{n}(j-j_n)^2
 \label{eq:gaussiangap}
\end{equation}
for a fixed \(c_0>0\) and all sufficiently large \(n\).
\end{lemma}
\begin{proof}
Stirling's formula gives, uniformly on any compact positive multiple of
\(n/\log n\),
\[
 T_{n,j}=\frac{e^{F_n(j)}}{\sqrt{2\pi(n-j)}}(1+O(n^{-1})).
\]
On such an interval, \(-F_n''(j)\asymp(\log n)^2/n\), proving
\eqref{eq:gaussiangap}. The third derivative is
\(O((\log n)^3/n^2)\), so on \(j=j_n+\sigma_n y\), with bounded \(y\),
\[
 F_n(j)=F_n(j_n)-y^2/2+O(n^{-1/2}|y|^3).
\]
Moreover, \((n-j)/j=r_n+O((\log n)/\sqrt n)\) on every bounded
standardized interval. The Gaussian Riemann sum, whose mesh is
\(1/\sigma_n\to0\), therefore contributes
\[
 \frac{e^{F_n(j_n)+cr_n}}{\sqrt{2\pi k_n}}
 \sqrt{2\pi}\sigma_n
 =S_ne^{cr_n}.
\]
The quadratic bound dominates larger standardized intervals within the
fixed relative neighborhood. On any other compact multiple of
\(n/\log n\), the scaled phase has a strict maximum at \(2\): with
\(y=j\log n/n\),
\[
 \frac{F_n(j)}n=\log n-2\log\log n+\log a+1+2\log y-y+o(1).
\]
This gives a fixed exponential loss outside a neighborhood of \(y=2\).
For \(j/n\) bounded below, the leading \(n\log n\) coefficient is smaller;
for smaller \(j\), the same expression and the limit \(2\log y-y\to-\infty\)
exclude the remaining range. These bounds also absorb the fixed weight
\(e^{c(n-j)/j}\) away from the saddle: for very small \(j\) it contributes
at most \(e^{|c|n}\), whereas the main logarithm is \(n\log n-O(n\log\log n)\).
Strict concavity handles the intermediate ranges. This proves the claims.
\end{proof}

\begin{remark}
The determinant factor \(D_n^{-1/2}\) is the product of the factorial
prefactor and the width of the discrete action saddle. Omitting either
factor changes the multiplicative constant. Replacing \(r_n\) in the large
exponential by \(\tfrac12\log n\) does not preserve an equivalent.
\end{remark}

\section{Absolute localization before cancellation}
\label{sec:localization}

\begin{theorem}[Arbitrarily accurate fixed-core localization]
\label{thm:localization}
Under \eqref{eq:assumptions}, there is a constant \(C>0\) such that
\begin{equation}
 A_n\le S_n n^C
 \qquad\text{for all sufficiently large }n.
 \label{eq:absolute-total}
\end{equation}
For every \(A>0\), there is a fixed integer \(M\ge J_0\) such that
\begin{equation}
 A_{n,M}^{(2)}\le S_n n^{-A}
 \qquad\text{for all sufficiently large }n.
 \label{eq:two-tail}
\end{equation}
The integer \(M\) may depend on the model and on \(A\), but not on \(n\).
\end{theorem}

The proof is given in several steps to make the absolute error scale
explicit. Constants below depend on the fixed model unless stated otherwise.
They never depend on the configuration or on \(n\). Constants carrying a
subscript \(M\) are permitted to depend on the fixed cutoff.

\subsection{A row bound and macroscopic localization}
Write the excesses of the nonunit actions as \(\ell_i=j_i-1\ge1\),
so \(L=\sum_i\ell_i\). The eventual polynomial hypothesis implies
\begin{equation}
 |E(\bm m)|\le a\sum_i\ell_i^2+C_1L+C_1
                  \le aL^2+C_1L+C_1.
 \label{eq:Ebound}
\end{equation}
Choose \(W\ge1\) bounding all weights. For a fixed \(L\) and number \(K\)
of nonunit actions,
\[
 \sum_{\substack{\bm m:\sum m_j=K\\\sum(j-1)m_j=L}}
 \frac1{\prod m_j!}=\frac1{K!}\binom{L-1}{K-1}.
\]
Consequently the total combinatorial factor is at most
\begin{align}
 \sum_{K=1}^L\frac{(L+K)!}{(L+1)!K!}
       W^K\binom{L-1}{K-1}
 &\le (8W)^L. \label{eq:entropy}
\end{align}
For example, bound \(\binom{L+K}{K}\le4^L\),
\(\binom{L-1}{K-1}\le2^L\), and use \(L/(L+1)\le1\).
It follows that
\begin{equation}
 A_{n,L}\le (8W)^L
 \frac{(aL^2+C_1L+C_1)^{n-1-L}}{(n-1-L)!}.
 \label{eq:row}
\end{equation}

\begin{lemma}[Macroscopic absolute localization]
\label{lem:macro}
Fix sufficiently small \(\delta,\varepsilon>0\). Except for absolute mass
at most \(S_n e^{-c n}\), with some \(c>0\), every configuration satisfies
\begin{equation}
 |L+1-j_n|\le\delta j_n,\qquad
 \ell_{\max}\ge(1-\varepsilon)L.
 \label{eq:macro}
\end{equation}
The largest excess is unique when \(\varepsilon<1/2\).
\end{lemma}
\begin{proof}
We detail the localization implicit in \eqref{eq:row}. A row with
\(L\ge\eta n\), for fixed \(\eta>0\), has logarithm at most
\((1-\eta)n\log n+O(n\log\log n+n)\), so cannot approach the bare maximum.
For rows within \(O(n)\) of that maximum, necessarily \(L/n\to0\) and
\(\log L=\log n-O(\log\log n)\): otherwise the base
\(aL^2+O(L+1)\), divided by a factorial with index asymptotic to \(n\),
loses a positive multiple of \(n\log\log n\) or \(n\log n\).
In this restricted range, logarithmic expansion of the base and Stirling's
formula give the same leading variational expression as in
Lemma~\ref{lem:bare}, with \(y=L\log n/n\):
\begin{equation}
 \frac1n\log A_{n,L}
 \le\log n-2\log\log n+\log a+1+2\log y-y+o(1).
 \label{eq:macro-phase}
\end{equation}
Here the entropy \(O(L)\), the linear correction in the base, and the
factorial error are all \(o(n)\). One can first force
\(y=O(\log\log n)\) by comparison with the maximum, so the expansion is
valid along every potentially maximizing sequence; then
\(2\log y-y\to-\infty\) at both ends forces \(y\) into a compact
subset of \((0,\infty)\). Its unique strict maximum is at \(y=2\).
This proves a fixed exponential loss outside the first condition of
\eqref{eq:macro}; \(j_n\log n/n\to2\). Summing the at most \(n\) rows
only contributes \(O(\log n)\) to their logarithm.

If \(\ell_{\max}\le(1-\varepsilon)L\), convexity, or successively
transferring excess between two nonextreme parts, gives
\[
 \sum_i\ell_i^2\le\bigl((1-\varepsilon)^2+\varepsilon^2\bigr)L^2.
\]
Integer rounding changes this by \(O(L)\), which is harmless. The factor
in parentheses is strictly less than one. Replacing \(a\) by this factor
times \(a\) in \eqref{eq:macro-phase} lowers the maximum by a fixed
positive multiple of \(n\). This proves the second assertion and the
combined exponential estimate. Uniqueness follows because two excesses
larger than \((1-\varepsilon)L\) cannot sum to at most \(L\).
\end{proof}

\subsection{The cost of a background cloud}
In the region \eqref{eq:macro}, write
\[
 \ell_{\max}=L-B,\qquad
 B=\ell_1+\cdots+\ell_s\le\varepsilon L,
\]
where \(s=K-1\) counts the remaining actions. The largest action is in the
unperturbed tail for all sufficiently large \(n\). Fuse the entire
background into it, obtaining a single action \(L+1\). Its exponent is
\begin{equation}
 E_*(L)=\lambda_{L+1}-(L+2)\lambda_1
       =a(L+1)^2+(b-\lambda_1)(L+1)+d-\lambda_1.
 \label{eq:Estar}
\end{equation}
This quantity is positive and asymptotic to \(aL^2\). The difference in
exponents obeys
\begin{equation}
 E_*(L)-E(\bm m)\ge2aB(L-B)-C_2B.
 \label{eq:fusion}
\end{equation}
To check the finite-prefix issue explicitly, write
\(\lambda_j=aj^2+bj+d+\eta_j\), where \(\eta_j=0\) outside the prefix.
The quadratic part of the difference is
\(a(L^2-\sum_i\ell_i^2)\), and every other term in the difference is
bounded in absolute value by a constant times \(s\le B\).
Since \(\sum_{i=1}^s\ell_i^2\le B^2\), \eqref{eq:fusion} follows.
In this region \(E(\bm m)>0\) as well.

For \(h=n-1-L\), \eqref{eq:fusion} implies
\begin{equation}
 \left(\frac{E(\bm m)}{E_*(L)}\right)^h
 \le \exp\!\left[-\left(2-\frac{2B}{L}-\frac{C_3}{L}\right)
                    \frac{hB}{L}\right].
 \label{eq:cloudloss}
\end{equation}
The combinatorial coefficient relative to the single fused action is
\[
 \frac{(L+s+1)!}{(L+1)!\prod_{j\ne j_{\max}}m_j!}.
\]
After replacing multiplicities by ordered background excesses and dividing
by \(s!\), it is at most \((2WL)^s/s!\), including weights, because
\(s\le B\le\varepsilon L\). Thus, whenever
\(2-2B/L-C_3/L\ge\kappa>0\), the total relative background mass is at most
\begin{equation}
 \sum_{s\ge0}\frac{(2WL)^s}{s!}
 \left(\sum_{\ell\ge1}e^{-\kappa h\ell/L}\right)^s
 =\exp\!\left(\frac{2WL}{e^{\kappa h/L}-1}\right).
 \label{eq:cloudmajorant}
\end{equation}
Restrictions such as a prescribed total \(B\) or an upper bound on each
background excess were dropped only to increase this positive bound.

\subsection{A finer excess window}
Take, for definiteness, fixed sufficiently small
\(\delta,\varepsilon\le1/100\) in Lemma~\ref{lem:macro}.
For all sufficiently large \(n\) we can use \(\kappa=2-4\varepsilon\)
in \eqref{eq:cloudmajorant}. Uniformly in the macroscopic excess window,
\begin{equation}
 \frac{2WL}{e^{\kappa h/L}-1}\le n^{1/3}.
 \label{eq:roughcloud}
\end{equation}
Indeed, \(L\asymp n/\log n\) and
\(h/L\ge r_n/(1+\delta)+O(1)\); the power of \(n\) on the left is
strictly smaller than \(1/3\) for these choices.
Moreover
\begin{equation}
 \frac{E_*(L)^h}{h!}
   =T_{n,L+1}\exp(O(\log n))
 \label{eq:reference-row}
\end{equation}
uniformly in that window. By \eqref{eq:gaussiangap}, configurations outside
\begin{equation}
 |L+1-j_n|\le\frac{j_n}{(\log n)^2}
 \label{eq:finewindow}
\end{equation}
but within \eqref{eq:macro} therefore have total mass at most
\begin{equation}
 S_n\exp\!\left(-c_1\frac{n}{(\log n)^4}+n^{1/3}+O(\log n)\right)
 \le S_n\exp\!\left(-c_2\frac{n}{(\log n)^4}\right).
 \label{eq:fineexception}
\end{equation}
This refinement is why an \(e^{o(n)}\) majorant alone is not enough: we need
to reach a window in which the remaining cloud costs only a power of \(n\).

\subsection{Excluding a mesoscopic background}
Inside \eqref{eq:finewindow}, consider \(B\ge L/\log n\), still with
\(B\le\varepsilon L\). Split the exponent in the coarse cloud bound in half.
The resulting total background mass, relative to \(E_*(L)^h/h!\), is at most
\begin{equation}
 \exp\!\left(-\frac{\kappa h}{2\log n}\right)
 \exp\!\left(\frac{2WL}{e^{\kappa h/(2L)}-1}\right)
 \le\exp\!\left(-c_3\frac{n}{\log n}\right).
 \label{eq:mesoscopic}
\end{equation}
For the last inequality, the positive logarithm is
\(O(n^{1-\kappa/4+o(1)})=o(n/\log n)\), whereas \(h\sim n\).
This is a direct exponential-moment bound, not a probabilistic independence
assumption.

\subsection{The polynomial cloud and marked-tail estimate}
It remains to consider \eqref{eq:finewindow} and \(B<L/\log n\).
Now \(h/L=r_n+O(1/\log n)\), and \eqref{eq:cloudloss} permits
\[
 z_n(L)=\exp\!\left[-\left(2-\frac2{\log n}-\frac{C_3}{L}\right)h/L\right]
 \le C_4e^{-2r_n}.
\]
The cloud exponent satisfies
\begin{equation}
 \mu_n(L)=\frac{2WLz_n(L)}{1-z_n(L)}\le C_5r_n,
 \label{eq:smallcloud}
\end{equation}
since
\(e^{-2r_n}=r_n(1+r_n)/(an)\) and \(L\sim n/(1+r_n)\).
Together with \eqref{eq:reference-row} and Lemma~\ref{lem:bare}, this proves
\eqref{eq:absolute-total}, including the previously discarded exceptions.

If there is a second primitive action \(j>M\), one background excess is at
least \(M\). Marking one such occurrence in the exponential majorant bounds
the relative mass by
\begin{equation}
 e^{\mu_n(L)}\frac{2WLz_n(L)^M}{1-z_n(L)}.
 \label{eq:markedcloud}
\end{equation}
The marked count is at least one on the event being bounded, so multiple
marks only overcount. Combining \eqref{eq:markedcloud} with the reference
row estimate yields, for every fixed \(M\ge J_0\),
\begin{equation}
 A_{n,M}^{(2)}\le
 C_M S_n\,r_n\exp\!\left(-(2M-2-C_6)r_n\right)
 +S_n\exp\!\left(-c_4\frac{n}{(\log n)^4}\right).
 \label{eq:tail-final}
\end{equation}
Here \(C_6\) is independent of \(M\); factors such as \(C_4^M\) are absorbed
in \(C_M\). The macroscopic exceptional set includes configurations with
no unique large action, so none have been omitted. Because
\(r_n\sim\tfrac12\log n\), taking \(M\) sufficiently large proves
\eqref{eq:two-tail} for any prescribed \(A\). This completes the proof of
Theorem~\ref{thm:localization}.\hfill\(\square\)

\begin{remark}[What this theorem does not optimize]
The proof gives existence of a sufficient fixed cutoff and arbitrary
polynomial accuracy. It does not identify the smallest sufficient cutoff.
In particular, it should not be confused with the sharp positive-coefficient
cutoff in the repository's forward theorem. Improving \(C_6\) or proving
that a much smaller inverse core always suffices remains a separate problem.
\end{remark}

\section{The analytic core sums the cancellations}
\label{sec:core}

Fix \(M\ge J_0\). Define
\begin{equation}
 F_M(q,u)=\sum_{j=1}^M w_jq^j e^{\lambda_j u},\qquad
 F_M(Q_M(u),u)=u,
 \qquad A_M(u)=\frac1{\partial_qF_M(Q_M(u),u)}.
 \label{eq:core}
\end{equation}
The analytic implicit-function theorem applies because
\(\partial_qF_M(0,0)=1\). Thus \(Q_M,A_M\) are analytic in a fixed
complex disc around zero. Their initial terms are
\begin{equation}
 Q_M(u)=u-\beta u^2+O(u^3),\qquad A_M(u)=1+O(u).
 \label{eq:corejets}
\end{equation}

\begin{lemma}[Exact one-tail decomposition]
\label{lem:decomposition}
Define the formal series
\begin{equation}
 V_M^{(1)}(u)=-A_M(u)\sum_{j>M}Q_M(u)^j e^{\lambda_j u}.
 \label{eq:onetail}
\end{equation}
Then
\begin{equation}
 \Q(u)=Q_M(u)+V_M^{(1)}(u)+R_M(u),
 \qquad |[u^n]R_M(u)|\le A_{n,M}^{(2)}.
 \label{eq:decomposition}
\end{equation}
Moreover, restricting the one-tail sum to
\(n/\log n\le j\le3n/\log n\) changes its degree-\(n\) coefficient by
\(O(S_ne^{-cn})\), for a suitable fixed \(c>0\).
\end{lemma}
\begin{proof}
Give every primitive action \(j>M\) a marker \(s\), and let \(Q_s\) be the
formal inverse defined by
\[
 u=F_M(Q_s,u)+s\sum_{j>M}Q_s^j e^{\lambda_j u}.
\]
Every fixed \(u\)-coefficient is a polynomial in \(s\). At \(s=0\), it is
\(Q_M\). Differentiating formally at zero gives \eqref{eq:onetail}.
Equivalently, in \eqref{eq:exact}, the marker degree counts occurrences of
actions above \(M\). The terms of marker degree at least two have absolute
mass bounded by \(A_{n,M}^{(2)}\), proving \eqref{eq:decomposition}.
If the sole tail action lies outside the stated interval, either \(L\) is
outside a fixed small neighborhood of \(j_n\), or the largest excess is
not within a fixed small relative distance of \(L\). For sufficiently large
\(n\), the tail action is the unique largest action on the complement of
these exceptional events: all other actions are at most the fixed \(M\).
Lemma~\ref{lem:macro} gives the final assertion.
\end{proof}

The expression \eqref{eq:onetail} is formal. We never integrate the entire
infinite tail on a positive circle, where it need not converge. Each
individual coefficient uses only \(j\le n\).

\subsection{A finite majorant for each core coefficient}
The next elementary coefficient estimate avoids a signed Cauchy-integral
argument and provides a useful certificate for an individual one-tail term.

\begin{lemma}[Analytic multiplier estimate]
\label{lem:multiplier}
Let \(A,H\) be analytic at zero, real on a real neighborhood, with
\(A(0)=H(0)=1\) and \(H'(0)=-\beta\). For a fixed analytic logarithm, set
\[
 B(u)=\frac{\log H(u)+\beta u}{u^2},
 \quad A^*(t)=\sum_{\ell\ge0}|[u^\ell]A(u)|t^\ell,
 \quad B^*(t)=\sum_{\ell\ge0}|[u^\ell]B(u)|t^\ell.
\]
Suppose \(L_j=aj^2+(b-\beta)j+d>0\), \(k\ge0\), and
\(t=k/L_j\) is inside the convergence radii of both majorants. Then
\begin{equation}
 \left|\frac{[u^k]A(u)H(u)^je^{(aj^2+bj+d)u}}
                   {L_j^k/k!}-1\right|
 \le A^*(t)e^{jt^2B^*(t)}-1.
 \label{eq:multiplierbound}
\end{equation}
In particular, uniformly for \(n/\log n\le j\le3n/\log n\), \(k=n-j\),
\begin{equation}
 [u^k]A(u)H(u)^je^{(aj^2+bj+d)u}
 =\frac{(aj^2)^k}{k!}
    e^{(b-\beta)k/(aj)}(1+O((\log n)^3/n)).
 \label{eq:multiplierasymp}
\end{equation}
\end{lemma}
\begin{proof}
Write the expression as \(e^{L_ju}G_j(u)\), where
\(G_j=A\exp(ju^2B)\) has constant term one. If \(g_{j,\ell}=[u^\ell]G_j\),
then the normalized coefficient is exactly
\[
 \sum_{\ell=0}^k g_{j,\ell}\frac{(k)_\ell}{L_j^\ell}.
\]
Here \((k)_\ell\le k^\ell\). The absolute coefficient majorant of
\(G_j\) at \(t\) is at most \(A^*(t)\exp(jt^2B^*(t))\), proving
\eqref{eq:multiplierbound}. In the stated range,
\(t=O((\log n)^2/n)\) and \(jt^2=O((\log n)^3/n)\).
Finally,
\[
 \left(\frac{L_j}{aj^2}\right)^k
 =\exp\!\left(\frac{b-\beta}{a}\frac{k}{j}
       +O(k/j^2)\right),\qquad k/j^2=O((\log n)^2/n).
\]
This proves \eqref{eq:multiplierasymp}. The bound also covers \(k=0\),
when \(t=0\) and the coefficient is exactly one.
\end{proof}

\subsection{Proof of the main theorem}
In Lemma~\ref{lem:multiplier}, use
\(A=A_M\) and \(H=Q_M/u\). Equations \eqref{eq:corejets} give its hypotheses.
For the finite central range of \(j\), the coefficient in \eqref{eq:onetail}
is therefore
\begin{equation}
 [u^n]V_M^{(1)}
 =-\sum_{n/\log n\le j\le3n/\log n}
 T_{n,j}e^{\gamma(n-j)/j}
       (1+O((\log n)^3/n))+O(S_ne^{-cn}).
 \label{eq:onetailasymp}
\end{equation}
The principal summands are positive before the common minus sign, and the
relative error is uniform. Lemma~\ref{lem:bare} gives
\begin{equation}
 [u^n]V_M^{(1)}\sim-S_ne^{\gamma r_n}.
 \label{eq:onetailresult}
\end{equation}
Choose \(A>|\gamma|/2+2\) in Theorem~\ref{thm:localization}, and then choose
its fixed \(M\). The remainder in \eqref{eq:decomposition} is
\(o(S_ne^{\gamma r_n})\), since \(r_n\sim\tfrac12\log n\).
The analytic coefficient \([u^n]Q_M\) is bounded by \(C_M^n\), while
\(S_n^{1/n}\to\infty\); it is negligible on the same scale.
Equation \eqref{eq:main} follows. Its right-hand side is strictly negative,
proving eventual negativity. Corollaries~\ref{cor:pure} and
\ref{cor:prefix} follow immediately.\hfill\(\square\)

\begin{remark}[Why the two exponential corrections differ]
The forward core inserts \(G(q)=q+\beta q^2+\cdots\) inside
\(e^{aj^2G(q)}\), generating a term of order \(\beta r_n^2/a\).
The inverse core inserts \(Q_M(u)=u(1-\beta u+\cdots)\) into a large
power \(Q_M(u)^j\), generating \(-\beta r_n/a\).
These are different mechanisms. A formal inverse identity, without a tail
estimate, does not license replacing one by the other.
\end{remark}

\section{Universal cancellation and what it does not imply}
\label{sec:cancellation}

In the pure model, combine Corollary~\ref{cor:pure} with the forward
quadratic equivalent proved in the source article \cite{finitecore}:
\begin{equation}
 u_n\sim S_n e^{(1+1/a)r_n^2}.
 \label{eq:forwardimport}
\end{equation}
This is the only consequence in the present report that imports a
nonclassical asymptotic theorem rather than proving its input anew.
It gives
\begin{equation}
 \boxed{\displaystyle
 -\frac{v_n}{u_n}\sim
 \exp\!\left(-(1+1/a)(r_n^2+r_n)\right).}
 \label{eq:cancellation}
\end{equation}
The factor tends to zero faster than every fixed power of \(n^{-1}\),
because \(r_n\sim\frac12\log n\). Nevertheless, its logarithm is \(o(n)\).
Thus the forward and inverse sequences can have identical exponential
coefficient type on a factorial--logarithmic scale while differing by a
factor invisible to that invariant. This is a concrete obstruction to
using type preservation as a substitute for sharp inversion.

There is also a hierarchy of finite-prefix effects. Altering \(\lambda_1\)
or \(w_2\) changes \(\beta\), and therefore changes the leading inverse
coefficient by a factor \(\exp(-\Delta\beta\,r_n/a)\). Altering any other
finite set of primitive data, while keeping \(\beta\) unchanged, changes the
coefficient only by a relative \(o(1)\). The core contains every exceptional
primitive datum, but only its quadratic jet survives at this precision.

The tail parameter \(b\) and the core parameter \(\beta\) occur through the
single combination \(b-\beta\). This is not an arbitrary normalization:
\(bj\,u\) from the tail and \(-\beta j\,u\) from \(Q_M(u)^j\) have the
same size at the coefficient saddle. The constant tail term \(du\), and
all higher core jets, vanish in the relative leading asymptotic.

None of these results asserts convergence of the original kernel at
positive arguments. If \(q,u>0\), the eventually quadratic summands
\(q^j e^{(aj^2+bj+d)u}\) fail even to tend to zero. All uses of the
countably infinite equation in the proof are formal and coefficientwise;
all analytic evaluations use a fixed finite core.

\section{The inverse Borel transform}
\label{sec:borel}

This section proves Theorem~\ref{thm:borel}. It also explains why an entire
Borel transform does not by itself provide a Borel sum in a specified
direction.

\subsection{Envelope identities}
For real \(x>0\), continue the saddle coordinate using \eqref{eq:r}, and
write
\[
 E(x)=\frac{x r(x)}{1+r(x)}(1+2r(x)),\qquad
 D(x)=1+4r(x)+2r(x)^2.
\]
Differentiating the defining equation of \(r\) gives
\begin{equation}
 r'(x)=\frac{r(x)(1+r(x))}{xD(x)}.
 \label{eq:rprime}
\end{equation}
The stationary-action identity in Lemma~\ref{lem:bare}, or direct
differentiation, then gives
\begin{equation}
 E'(x)=2r(x),\qquad
 E''(x)=\frac{2r(x)(1+r(x))}{xD(x)}>0.
 \label{eq:envelope}
\end{equation}
These identities preserve every slowly varying term needed below.

By Theorem~\ref{thm:main} and Stirling's formula,
\[
 \left(\frac{|v_n|}{n!}\right)^{1/n}
 \sim\frac{e^{2r_n+1/(1+r_n)}}{n}
 =\frac{a e^{1/(1+r_n)}}{r_n(1+r_n)}\longrightarrow0.
\]
This proves that \(\Bcal\Q\) is entire. In contrast,
\(|v_n|^{1/n}\sim e^{2r_n-1+1/(1+r_n)}\to\infty\), so the original
formal inverse has radius of convergence zero.

\subsection{A second, positive coefficient saddle}
For \(z>0\), define
\begin{equation}
 \Phi_z(x)=E(x)-x\log x+x+x\log z.
 \label{eq:borelphase}
\end{equation}
Equations \eqref{eq:rprime}--\eqref{eq:envelope} imply
\begin{align}
 \Phi_z'(x)
 &=2r(x)-\log x+\log z
   =\log\frac{az}{r(x)(1+r(x))},\\
 \Phi_z''(x)&=-\frac{2r(x)+1}{xD(x)}<0.
 \label{eq:borelderivatives}
\end{align}
Thus there is a unique real maximum. If \(R(R+1)=az\), it is at
\begin{equation}
 m=m(z)=z e^{2R}=\frac{R(R+1)}a e^{2R}.
 \label{eq:borelm}
\end{equation}
At this point,
\begin{equation}
 \Phi_z(m)=\frac Ra e^{2R},\qquad
 \sigma_B^2=\frac1{-\Phi_z''(m)}
           =\frac{m(1+4R+2R^2)}{2R+1}.
 \label{eq:borelsaddle}
\end{equation}
In particular, \(\sigma_B\to\infty\) and \(\sigma_B/m\to0\).

Apply Stirling's formula to the reference sum
\begin{equation}
 H(z)=\sum_{n\ge1}\frac{S_ne^{\gamma r_n}}{n!}z^n.
 \label{eq:referenceborel}
\end{equation}
Uniformly for \(n\) near \(m\), its summand is
\[
 \frac{e^{\gamma r_n}}{\sqrt{2\pi nD_n}}
 e^{\Phi_z(n)}(1+O(1/n)).
\]
The logarithmic prefactor varies by \(o(1)\) over a fixed number of
standard deviations. On \([m/2,2m]\), differentiation of
\eqref{eq:borelderivatives} gives
\(\Phi_z'''(x)=O(1/(m^2R))\). Consequently its standardized cubic
coefficient is \(O(\sqrt{R/m})\to0\).
For example, on
\(\abs{n-m}\le\sigma_B(m/R)^{1/12}\), the cubic remainder is uniformly
\(o(1)\). Taylor expansion and the Gaussian Riemann sum therefore give
\begin{align}
 H(z)&\sim
 \frac{e^{\gamma R}}{\sqrt{2\pi m(1+4R+2R^2)}}
 e^{\Phi_z(m)}\sqrt{2\pi\sigma_B^2}\notag\\
 &=\frac{e^{\gamma R}}{\sqrt{2R+1}}
   \exp\!\left(\frac Ra e^{2R}\right).
 \label{eq:Hborel}
\end{align}

For completeness, the omitted tails cannot change this equivalent.
Within \([m/2,2m]\), the second derivative is bounded above and below by
negative constant multiples of \(1/(mR)\), giving a Gaussian upper bound
outside the central window. At either endpoint the phase has already lost
at least \(cm/R\). Outside that interval, concavity bounds the phase by
the tangent at the endpoint. Its outward slope has absolute value at least
\(c/R\), so the right tail is bounded by a geometric sum. Polynomial
factors from \(e^{\gamma r_n}/\sqrt{nD_n}\) are absorbed by the same bound;
on the finite left tail their logarithm is \(O(\log m)\), negligible
beside \(m/R\). These observations justify the sum asymptotic without an
unproved interchange of a nonuniform expansion and an infinite sum.

Finally, for every \(\epsilon>0\), Theorem~\ref{thm:main} gives a fixed
\(N_\epsilon\) such that
\[
 (1-\epsilon)S_ne^{\gamma r_n}
 \le -v_n\le(1+\epsilon)S_ne^{\gamma r_n}
 \quad(n\ge N_\epsilon).
\]
The finite initial part of the Borel series is a polynomial in \(z\),
negligible compared with \eqref{eq:Hborel}. Positivity of the reference
summands now sandwiches \(-\Bcal\Q(z)\) between
\((1\pm\epsilon)H(z)\), up to negligible terms. Letting \(\epsilon\to0\)
proves \eqref{eq:borel}.

\subsection{The positive-direction obstruction}
For a positive Laplace parameter \(t\), the ordinary factorial-normalized
Borel--Laplace prescription would be
\[
 \frac1t\int_0^\infty e^{-z/t}\Bcal\Q(z)\dd z.
\]
The equivalent \eqref{eq:borel} shows that its integrand is eventually
negative and has logarithm of absolute value
\[
 \frac{R(z)}a e^{2R(z)}-\frac zt+O(R(z)+\log R(z)).
\]
Since \(R(z)\sim\sqrt{az}\), this tends to positive infinity. The
improper integral diverges to \(-\infty\). There is no oscillatory
cancellation on this ray. This proves the final assertion of
Theorem~\ref{thm:borel}; it says nothing against summation in another
direction, nor against a different acceleration procedure.

\section{The smallest formal inverse term}
\label{sec:least}

Fix \(0<t<1\) and put \(R=\tfrac12\log(1/t)\). The reference term
associated with Theorem~\ref{thm:main} is
\begin{equation}
 h_t(x)=\frac{e^{\gamma r(x)}}{\sqrt{D(x)}}
             \exp(E(x)+x\log t).
 \label{eq:leastreference}
\end{equation}
The large phase \(\Psi_t(x)=E(x)+x\log t\) is strictly convex by
\eqref{eq:envelope}; its derivative is \(2(r(x)-R)\). Its unique minimum
is therefore attained at
\begin{equation}
 m=m(t)=\frac{R(R+1)}{at},\qquad
 \Psi_t(m)=-\frac{R^2}{at}.
 \label{eq:leastm}
\end{equation}
Moreover
\begin{equation}
 \Psi_t''(m)=\frac{2R(R+1)}{m(1+4R+2R^2)}\sim\frac1m.
 \label{eq:leastcurvature}
\end{equation}
The derivative of the logarithmic prefactor in \eqref{eq:leastreference}
is \(O(1/m)\) near \(m\). Its inclusion shifts the continuous reference
minimum by only \(O(1)\); rounding to a neighboring integer changes the
minimum value by \(1+o(1)\). To see that this is the global minimum, note
that, for all sufficiently large \(x\), the second derivative of
\(\log h_t(x)\) is \(E''(x)+O(x^{-2})>0\). On any fixed bounded range
of positive indices the reference value is only a fixed power of \(t\),
whereas \(h_t(m)\) is exponentially smaller. It follows that
\begin{equation}
 \min_{n\ge N}h_t(n)
 \sim\frac{\exp(-R^2/(at)+\gamma R)}{\sqrt{1+4R+2R^2}}.
 \label{eq:referencemin}
\end{equation}

We must still justify replacing the reference coefficients by the exact
ones inside a minimum. For every fixed \(t>0\),
\(|v_n|t^n\to\infty\), because \(|v_n|^{1/n}\to\infty\); hence a
minimum over \(n\ge N\) exists. As \(t\downarrow0\), any minimizing
index tends to infinity: every fixed finite set of terms is polynomially
small in \(t\), but the term near \(m(t)\) is smaller than every fixed
power of \(t\). On all sufficiently large indices, the exact-to-reference
ratio lies uniformly in \([1-\epsilon,1+\epsilon]\). The same sandwich
argument used for the Borel sum now applies to the minimum. It proves
\eqref{eq:least}.

For the index assertion, on the scale \(n=m+y\sqrt m\), Taylor's formula
and \eqref{eq:leastcurvature} give, uniformly for bounded real \(y\),
\begin{equation}
 \log\frac{h_t(n)}{h_t(m)}=\frac{y^2}{2}+o(1).
 \label{eq:leastgaussian}
\end{equation}
For any fixed \(\delta>0\), eventual convexity gives a strictly positive
logarithmic gap outside \(\abs{n-m}<\delta\sqrt m\), first locally by
\eqref{eq:leastgaussian} and then globally by monotonicity on each side
of the reference minimizer. The uniform relative error in the exact
coefficients tends to zero, so it cannot close this fixed gap. Thus
\(\abs{n_{\min}-m}<\delta\sqrt m\) eventually. Since \(\delta\) is
arbitrary, \eqref{eq:minindex} follows.

\begin{remark}[Three distinctions about optimal truncation]
The smallest coefficient term is not automatically the error of a truncated
analytic function. Such a conclusion requires a specified analytic
realization and a remainder theorem. Also, an equivalent with unspecified
relative \(o(1)\) error does not determine the minimizing integer to
\(O(1)\), because the curvature is only of order \(1/m\). Finally,
finite initial zero coefficients would make an unrestricted minimum equal
to zero; this is why the theorem is stated after an eventual-negative index
\(N\). None of these issues is hidden by the notation ``least term.''
\end{remark}

% ed. (2026-09-29): relation to the negative-ray article's remainder bound, uncited by this article.
\begin{ednote}
At \(a=1\), \(b=d=0\) the exponent \(-\log^2(1/t)/(4t)\) of
\eqref{eq:least} has the constant \(1/4\) of the negative-ray article's
remainder bound \cite{ed:nrs} (its Theorem~\texttt{thm:optimal}), which is
an upper bound from a majorant. The present theorem is a statement about
formal terms and gives no lower bound for the analytic truncation error.
\end{ednote}

\section{Exact computations and reproducibility}
\label{sec:verification}

\subsection{An integer-normalized recurrence}
The program \path{code/verify.py} computes
\(V_n=n!v_n\) and \(P_{j,n}=n![u^n]\Q(u)^j\). Differentiating the
power gives the exact recurrence
\begin{equation}
 P_{j,n}=j\sum_{k=1}^{n-j+1}\binom{n-1}{k-1}V_kP_{j-1,n-k}
 \qquad(j\ge2),
 \label{eq:powerrecurrence}
\end{equation}
with \(P_{1,n}=V_n\), \(V_1=1\), and vanishing coefficients below the
natural degree. The inverse kernel then gives, for \(n\ge2\),
\begin{align}
 V_n={}&-\sum_{k=1}^{n-1}\binom nk V_k\lambda_1^{n-k}\notag\\
 &-\sum_{j=2}^n w_j\sum_{k=j}^n
               \binom nk P_{j,k}\lambda_j^{n-k}.
 \label{eq:integerrecurrence}
\end{align}
The right-hand side only contains previously available data. Integer
primitive parameters give integers throughout, so there is no roundoff
in the coefficient files. The algorithm uses polynomially many arithmetic
operations, although the cost of large-integer arithmetic grows with the
size of these integers. No unit-cost bit-complexity claim is intended.

\subsection{Independent finite checks}
The supplied run used five models: \(\lambda_j=j^2\);
\(\lambda_j=2j^2\); \(\lambda_1=0\) with \(\lambda_j=j^2\) for
\(j\ge2\); \(\lambda_j=j^2+j\); and \(w_2=2\) with all other
weights one and \(\lambda_j=j^2\). The main pure model was computed
through degree 320, and each additional model through degree 160.

There were 255 exact assertions in the recorded run. For each model, the
integer recurrence was compared through degree 14 with the independent
finite partition formula \eqref{eq:exact}. The marked one-tail
formula was separately checked through degree 12 at cutoffs \(M=1,2,3\)
against exact truncated-polynomial evaluation of the analytic-core response.
The small cutoff checks test the exact identity, not the asymptotic
sufficiency of those particular cutoffs. All assertions passed.
The check program does not serve as a machine proof of the asymptotic
localization theorem.

The files \path{data/*_exact_egf.csv} contain \(V_n=n!v_n\), not \(v_n\).
This normalization is recorded explicitly to prevent a factorial from
being mistaken for a disagreement in coefficient growth. The file
\path{data/verification.json} records the check count and run parameters.

\subsection{Asymptotic diagnostics}
Table~\ref{tab:pure} compares the exact pure-model coefficients with the
proved equivalent. Decimal evaluation uses 90-digit \texttt{mpmath}
arithmetic. These numbers are diagnostic approximations, not intervals.

\begin{table}[htbp]
\centering
\caption{Pure model \(a=1\): exact coefficients divided by the proposed
and now proved leading equivalent. The slow approach is visible.}
\label{tab:pure}
\begin{tabular}{rrr}
\toprule
\(n\)&\(r_n\)&\(v_n/(-S_ne^{-2r_n})\)\\
\midrule
20 & 1.087760 & 0.448613578 \\
50 & 1.368199 & 0.649597330 \\
100 & 1.593244 & 0.757155401 \\
160 & 1.751348 & 0.813338189 \\
200 & 1.827836 & 0.835909537 \\
240 & 1.890976 & 0.852577838 \\
300 & 1.969011 & 0.870968827 \\
320 & 1.991732 & 0.875899988 \\
\bottomrule
\end{tabular}
\end{table}

The agreement of the rows at \(n=50,100,200,300\) with the earlier
source's conjectural diagnostic table \cite{finitecore} is a useful
normalization check. The ratio at degree 320 is still about \(0.876\),
so the leading equivalent is not a high-accuracy finite-order approximation
at these degrees. The theorem claims a limit, not a certified onset value.

For the four other models, the ratios at degree 160 to
\(-S_ne^{(b-\beta)r_n/a}\) are respectively about
\(0.84676\), \(0.76341\), \(0.78906\), and \(0.84791\), in the order
listed after the main model above. These are useful tests of the core jet,
the slope coefficient, and the constant in the exponent; a finite table
cannot prove the universal perturbation theorem.

\subsection{Rebuilding the package}
From the package directory, run
\begin{quote}\ttfamily
python code/verify.py --degree 320\\
pdflatex -interaction=nonstopmode -halt-on-error article.tex\\
pdflatex -interaction=nonstopmode -halt-on-error article.tex
\end{quote}
A third LaTeX pass may be needed after changing the table of contents.
The exact checks use Python's standard-library integers and rational
arithmetic. Only the diagnostic decimal evaluations require
\texttt{mpmath}. The program performs no network requests and does not
require a local ProveIt checkout.

\section{Proof audit and a formalization route}
\label{sec:audit}

\subsection{Dependencies and scope}
The main coefficient theorem depends on the exact formal configuration
identity, elementary analytic implicit inversion for a finite core, an
absolute background-cloud bound, and a one-dimensional discrete saddle
estimate. No forward coefficient asymptotic or conjectural source theorem
is used in its proof. The two Borel and least-term theorems depend only on
the proved inverse equivalent and elementary saddle analysis.

The forward-to-inverse ratio \eqref{eq:cancellation} additionally uses
\eqref{eq:forwardimport} from the repository. Its provenance is separated
because proving the inverse conjecture does not retroactively formalize or
independently audit every forward result in its source.

The proof's most delicate step is Theorem~\ref{thm:localization}. Its
successive bounds are stated on progressively smaller sets: macroscopic
excess and a unique large action, a fine excess window, a submesoscopic
background, and finally a polynomial cloud. At each removal, the mass is
estimated before cancellation. The last estimate retains a cutoff-dependent
factor \(e^{-2Mr_n}\); the positive cloud constant does not depend on
\(M\). This independence is what permits an arbitrarily accurate fixed
cutoff. An \(e^{o(n)}\) estimate alone would not suffice.

No all-order asymptotic expansion, numerical threshold for the asymptotic
regime, smallest sufficient inverse core, or open-sector Borel continuation
is claimed. The proofs are conventional mathematical proofs supplied for
independent checking; the report contains no Lean code and no claim of a
completed formal verification.

\subsection{A modular route to Lean}
A practical formalization should not start by encoding an infinite saddle
sum all at once. The finite algebraic portion can be separated completely
from the asymptotics.

First, formalize the coefficientwise recursion for \(\U\), its compositional
inverse, and the marker derivative of a finite implicit kernel. The
recurrence \eqref{eq:integerrecurrence} and finite partition formula give
two executable representations that can be compared at arbitrary finite
truncation orders. Formalizing their equality is stronger than checking
selected coefficients.

Second, state the majorant theorem for finite weighted configurations of
a fixed total excess. Its essential ingredients are the fusion inequality,
the factorial-ratio bound, and a marked exponential generating-function
inequality. This finite theorem can expose all constants and avoid Hahn
series entirely. The countably infinite model is recovered coefficient by
coefficient.

Third, formalize a reusable discrete Laplace lemma for a strictly concave
phase with a diverging Gaussian width, uniform local Taylor bounds, and
explicit tail inequalities. The action saddle and Borel saddle then become
two instances with different curvature scales. A corresponding convex
minimum lemma handles the least term.

Finally, connect the coefficient field to the repository's transseries
support API by substituting \(q=e^{-x}\). This last bridge should state
only the formal series object and its coefficient theorems. An analytic
realization requires its own theorem and should not be silently attached
to the formal representation.

\section{Further research questions}
\label{sec:questions}

The questions below are continuations not resolved by this report. The first
three are the most direct extensions of the proof; the later ones connect
sharp coefficients to summation, effective computation, and transseries
algebra.

\begin{question}[The smallest sufficient inverse core]
For the pure quadratic model, is a two-action core already sufficient for a
full multiplicative inverse equivalent? More generally, characterize the
smallest \(M\) for which the first marked response has the same leading
coefficient as the inverse. The proof here chooses a possibly large cutoff
to overcome a deliberately coarse absolute majorant. That choice is an
existence result, not a sharp threshold. A smaller cutoff likely requires
retaining cancellations between additional tails instead of estimating
those tails absolutely.
\end{question}
% ed. (2026-09-29): this question is answered by two earlier-filed packages (batch 49).
\begin{ednote}
Answered in two earlier-filed packages. \cite{ed:sqi},
Theorem~\texttt{thm:two}: every fixed \(M\ge2\) gives the full inverse
equivalent, \(M=2\) is the smallest sufficient core in the pure model, and
\(M=1\) suffices exactly when \(w_2=0\). \cite{ed:qfr},
Theorem~\texttt{thm:sharp} and Corollary~\texttt{cor:minimalcore}: the
sharp relative error of every fixed core with \(M\ge2\), in the pure model
\(j_nt_n^M\).
\end{ednote}

\begin{question}[A complete inverse large-order expansion]
Determine an all-order expansion of \(v_n/(-S_ne^{\gamma r_n})\) in
explicit functions of \(r_n\) and \(n^{-1}\). The next terms receive
contributions from the cubic jet of the inverse core, the constant tail
term \(d\), the Stirling correction, and the displacement of the action
saddle caused by \(e^{\gamma(n-j)/j}\). Their relative sizes must be
ordered before truncation. A useful initial target is a first correction
that explains the visible error in Table~\ref{tab:pure}, with a proved
uniform remainder and no empirically fitted constants.
\end{question}

\begin{question}[Subquadratic signed condensation]
For \(\lambda_j=aj^p\), \(1<p<2\), develop the analogue of the absolute
localization theorem at a scale fine enough to survive cancellation.
At the bare saddle, the inverse-core perturbation \(j\log(Q_M(u)/u)\)
can contain several terms that do not tend to zero. Identify the necessary
core jet order and determine the resulting cancellation exponent. The
forward finite-core threshold cannot simply be copied: forward and inverse
cores enter different large powers.
\end{question}

\begin{question}[Slowly varying departures from a quadratic tail]
Replace the eventual polynomial identity by
\(\lambda_j=aj^2+bj+\eta j/\log j+o(j/\log j)\), with suitable control
of local oscillations. Since \((n-j)/(j\log j)\to1/2\) at the bare
saddle, the direct one-tail calculation predicts an extra factor
\(e^{\eta/(2a)}\). Prove a uniform fusion bound for this class and determine
which finite-difference assumptions are actually necessary. More irregular
perturbations could destroy a smooth local saddle while preserving a
logarithmic coefficient law.
\end{question}
% ed. (2026-09-29): this question is answered by an earlier-filed package (batch 49).
\begin{ednote}
Answered in \cite{ed:qfr}, Theorem~\texttt{thm:borderline}: if
\(\lambda_j=aj^2+bj+d+\varepsilon_j\) (in this article's notation) with
\(\varepsilon_j\log j/j\to\kappa\), with no monotonicity or
differentiability assumption on \(\varepsilon_j\), the inverse equivalent
acquires exactly the factor \(e^{\kappa/(2a)}\), and so does the Borel
equivalent.
\end{ednote}

\begin{question}[Effective, rather than asymptotic, cancellation bounds]
Make the constants in \eqref{eq:fusion}, \eqref{eq:cloudloss}, and
\eqref{eq:tail-final} explicit for the pure model. Combine them with
interval bounds for the saddle sum and the analytic core to produce a
computable enclosure for \(v_n/(-S_ne^{-2r_n})\). The target is an
inequality valid for a stated finite range of degrees, not a floating-point
agreement or a nonconstructive ``sufficiently large'' statement.
\end{question}

\begin{question}[Angular inverse Borel growth]
Theorem~\ref{thm:borel} concerns one real ray. Determine the growth of
\(\Bcal\Q(z)\) for complex \(z\) in sectors, including the transition
between constructive and destructive phase interference. Eventual signs
reduce the positive ray to a real saddle but provide no lower bound off
that ray. A solution could locate directions of ordinary summability or
identify a suitable acceleration scale without confusing entire
continuation with acceptable growth.
\end{question}
% ed. (2026-09-29): partial answers in two earlier-filed packages (batch 49).
\begin{ednote}
Partly answered. \cite{ed:sqi}, Theorem~\texttt{thm:borel}, shows (for an
exact tail \(aj^2\)) that the maximum modulus of \(\Bcal\Q\) on the circle
\(|z|=R\) is asymptotic to \(-\Bcal\Q(R)\). For the pure model at \(a=1\),
\cite{ed:nbq}, Theorem~\texttt{thm:main}(iii), shows that \(\Bcal\Q\) (in
the normalization \eqref{eq:borel-def}) admits no exponential bound in any
sector of positive opening about the direction \(\pi\). Growth in general
sectors remains open.
\end{ednote}

\begin{question}[Least term versus an actual analytic remainder]
For a specified negative-direction analytic realization of the inverse,
determine whether the scale in \eqref{eq:least} is attained by its optimal
truncation error, or whether the error has an additional factor. This
requires both an upper remainder estimate and a matching lower bound.
The coefficient theorem alone gives neither. Any comparison should specify
the sector, the branch, and whether the truncation index is chosen
continuously or integrally.
\end{question}

\begin{question}[A moving-saddle asymptotic algebra]
Construct an algebra of coefficient asymptotics that retains an implicit
slowly varying coordinate such as \(r(1+r)e^{2r}=an\), rather than only a
fixed factorial scale. Determine composition, inversion, differentiation,
and multiplication rules that preserve enough information to recover
\eqref{eq:cancellation}. A useful goal is an effective analogue of the
factorial-asymptotic calculus in \cite{borinsky}, but with a nontrivial
invariant for the zero-factorial-type sequences considered here.
\end{question}

\begin{question}[Robustness under signed or oscillatory primitive data]
How far can the assumptions \(w_j,\lambda_j\ge0\) be relaxed? Finitely
many signed primitive weights may still admit a controlled absolute
majorant, but the quadratic core jet can vanish or change sign, and
oscillatory tail weights may introduce several equally dominant saddles.
Classify the cases with eventual inverse sign, periodic sign patterns, or
persistent cancellation at leading order. This question requires a new
statement rather than a formal replacement of positive constants by
absolute values.
\end{question}
% ed. (2026-09-29): partial answer in an earlier-filed package (batch 49).
\begin{ednote}
Partly answered in \cite{ed:sqi}, Theorem~\texttt{thm:main}: with an exact
tail \(w_j=1\), \(\lambda_j=aj^2\), finitely many exceptional weights and
slopes may be arbitrary real numbers, and the inverse coefficients are still
eventually negative, with \(v_n\sim-S_n\exp(-(\lambda_1+w_2)r_n/a)\).
Oscillatory tail weights remain open.
\end{ednote}

\section{Conclusion}
The sharp quadratic inverse conjecture is reduced to, and proved by, an
absolute localization estimate strong enough to survive the signed
cancellation. A fixed finite core incorporates every exceptional primitive
parameter, yet only its quadratic coefficient survives in the leading
inverse law. The same law determines an entire Borel transform whose
positive-axis growth is too large for ordinary Laplace summation, and a
precise smallest-formal-term scale.

The principal mathematical distinction is between preserving a coefficient
class and preserving a coefficient equivalent. In this model the difference
is a factor of size \(\exp(-c(\log n)^2)\) relative to the forward
series. Capturing that factor requires summing the analytic-core
cancellations and bounding the omitted signed configurations on the
smaller scale, not merely transferring a Gevrey estimate.

\begin{thebibliography}{9}
\raggedright
\addcontentsline{toc}{section}{References}

\bibitem{repo}
V. Reshetnikov, \emph{ProveIt}, GitHub repository.
Snapshot \texttt{\RepoCommit}, inspected September 29, 2026.
\url{https://github.com/VladimirReshetnikov/ProveIt}.

\bibitem{inventory}
\emph{ProveIt transseries research inventory},
\path{Analysis/Transseries/docs/series-and-transseries/README.md},
same snapshot as \cite{repo}.
\href{https://github.com/VladimirReshetnikov/ProveIt/blob/aab173a7ccdd48add72191d68ce6b0c12ffd72fe/Analysis/Transseries/docs/series-and-transseries/README.md}{Pinned source}.
In particular, the descriptions of the finite-core, microscopic-condensation,
negative-ray-summation, and sharp-weighted-type research packages.

\bibitem{finitecore}
\emph{Finite-Core Universality and Sharp Large Order for Countable
Exponential-Feedback Transseries}, ProveIt research manuscript, 2026,
\path{Finite_Core_Universality_Exponential_Feedback/article.tex},
under the directory in \cite{inventory}, same pinned snapshot.
\href{https://github.com/VladimirReshetnikov/ProveIt/blob/aab173a7ccdd48add72191d68ce6b0c12ffd72fe/Analysis/Transseries/docs/series-and-transseries/Finite_Core_Universality_Exponential_Feedback/article.tex}{Pinned source}.
Theorem labeled \texttt{thm:explicit}; proposition
\texttt{prop:inverse-response}; conjecture \texttt{conj:quadratic-inverse}.

% ed. (2026-09-29): bibliography entries for the editorial notes (packages filed beside this one).
\bibitem{ed:qfr}
\emph{Quadratic Exponential Feedback after Reversion}, ProveIt research
package filed in batch 49 (2026-09-29),
\path{Analysis/Transseries/docs/series-and-transseries/Quadratic_Exponential_Feedback_After_Reversion/article.tex}.
Editorial addition.

\bibitem{ed:sqi}
\emph{Signed Quadratic-Feedback Inversion}, ProveIt research package filed
in batch 49 (2026-09-29),
\path{Analysis/Transseries/docs/series-and-transseries/Signed_Quadratic_Feedback_Inversion/article.tex}.
Editorial addition.

\bibitem{ed:nbq}
\emph{Natural Boundaries of Quadratic Exponential Feedback}, ProveIt research
package filed in batch 49 (2026-09-29),
\path{Analysis/Transseries/docs/series-and-transseries/Natural_Boundaries_Quadratic_Exponential_Feedback/article.tex}.
Editorial addition.

\bibitem{ed:nrs}
\emph{Negative-Ray Summation of Countable Exponential-Feedback Transseries}, ProveIt
research package filed in batch 47 (2026-09-29),
\path{Analysis/Transseries/docs/series-and-transseries/Negative_Ray_Summation_Exponential_Feedback/article.tex}.
Editorial addition.

\bibitem{gessel}
I. M. Gessel, \emph{Lagrange inversion},
Journal of Combinatorial Theory, Series A \textbf{144} (2016), 212--249.
\href{https://doi.org/10.1016/j.jcta.2016.06.018}{doi:10.1016/j.jcta.2016.06.018}.
\url{https://arxiv.org/abs/1609.05988}.

\bibitem{jkt}
S. Jansen, T. Kuna, and D. Tsagkarogiannis,
\emph{Lagrange inversion and combinatorial species with uncountable color
palette}, Annales Henri Poincar\'e (2021).
\href{https://doi.org/10.1007/s00023-020-01013-0}{doi:10.1007/s00023-020-01013-0}.
\url{https://arxiv.org/abs/2008.10862}.

\bibitem{borinsky}
M. Borinsky, \emph{Generating asymptotics for factorially divergent
sequences}, Electronic Journal of Combinatorics \textbf{25}(4) (2018), P4.1.
\href{https://doi.org/10.37236/5999}{doi:10.37236/5999}.
\url{https://arxiv.org/abs/1603.01236}.

\bibitem{edgar}
G. A. Edgar, \emph{Transseries for beginners},
Real Analysis Exchange \textbf{35} (2010), 253--310.
\url{https://arxiv.org/abs/0801.4877}.

\end{thebibliography}
\end{document}
